% Level 1, Week 12 -- Showcase and Celebration   (input by lesson-plans.tex)
\lesson{12}{Showcase and Celebration}{Share what you made, explain how it works, and look back at how far you have come.}
\lessoninfo{Concepts: explaining code (R), reflection, exit assessment}{Tools: Thonny + turtle, the students' projects}{CSTA: 1B-AP-17, 1B-IC-20, 1B-NI-05}

\begin{objectives}
  \item The student can run their project for a visitor and explain one part of the code in their own words (Milestone M4) (L1.G7, L1.G9; L1.W12.1)
  \item The student can give and receive \enquote{two stars and a wish} feedback during a gallery walk (L1.G7, L1.G8; L1.W12.2)
  \item The student can read, predict and fix short Python programs in the exit assessment (L1.G1--G6, L1.G9; L1.W12.3)
  \item The student can share work safely: on screen only, nicknames only, no personal information (L1.G8; L1.W12.4)
\end{objectives}

\begin{materials}
  \item Print: the \textbf{Exit assessment in the Assessments booklet} (one per student; answers and rubric in the Answer Key), feedback cards (last page of the Project Teacher Guide, 4 per student), certificates (\texttt{extras/certificate}, write names by hand), workbook chapter 12 and homework 12.
  \item One \enquote{stand} per project: a Mac with Thonny open and the project loaded, the plan sheet beside it, and a folded name card with the project name and the artist's \textbf{nickname}.
  \item Optional: invite family members for minutes 10--35 (the gallery). Snacks help!
  \item Read \enquote{Showcase plan} in the Project Teacher Guide. Prep time: about 10 minutes (plus printing).
\end{materials}

\needspace{20\baselineskip}
\section{At a glance}
Times are a guide. If families can only come at the end, do the exit assessment first (straight after the warm-up).
\begin{timeplan}
  \timeslot{0--7}{Warm-up}{Set up stands; last polish (no new features); practise the 2-minute demo in pairs.}
  \timeslot{7--10}{Concept}{How a gallery works; feedback cards; showcase safety.}
  \timeslot{10--13}{Guided demo}{You model a 2-minute demo with the Night Sky Postcard.}
  \timeslot{13--35}{Hands-on build: showcase}{Gallery walk: each artist demos to visitors; visitors fill in feedback cards.}
  \timeslot{35--55}{Exit assessment}{Students only, quiet, about 20 minutes (Assessments booklet).}
  \timeslot{55--60}{Wrap-up}{Read feedback, certificates, one thing I learned / want to make next, Level 2 preview, homework 12.}
\end{timeplan}

\section{Warm-up (0--7 min)}
\begin{freebuild}[Last polish (3 minutes, no new features!)]
Each artist runs their project once from start to finish, with a normal answer. Allowed: fixing a typo, changing a colour, adding a title. Not allowed: anything new and big --- \enquote{a showcase is not the time to break things}. Save (\keys{Cmd}+\keys{S}).
\end{freebuild}
\begin{pairroles}[Artist and assistant]
In pairs, practise the 2-minute demo script (in the workbook) once each: the \textbf{artist} talks and points, the \textbf{assistant} types the visitor's answer and presses Run. Swap. With 1 student: you are the assistant and the first visitor. With 3: artist, assistant and practice visitor rotate. With 4: two pairs.
\end{pairroles}

\section{Concept (7--10 min)}
\begin{teacherscript}
\enquote{Today you are artists at your own art show. Every visitor will try your program, and you'll show them one part of your code. Visitors write feedback cards: two stars --- two things they liked --- and a wish. When you are a visitor, be the kind of visitor you'd like to have.}
\end{teacherscript}
\begin{analogy}[A museum guide]
A museum guide doesn't explain every brush stroke. They pick one interesting thing and tell its story. Pick \textbf{one} part of your code --- your function, your loop or your \texttt{if} --- and tell its story.
\end{analogy}
\begin{safety}[Showcase safety]
Read aloud before guests arrive:
\begin{itemize}[leftmargin=1.2em,itemsep=1pt]
  \item Projects are shown \textbf{on the screen only}: nothing is uploaded, posted or emailed.
  \item Name cards and screens show \textbf{nicknames only}; programs never ask visitors for private information.
  \item Photos of the artwork only if the artist agrees, with no faces or names on screen, and they stay in the family.
  \item If you borrowed an idea, say so in your demo: \enquote{The house idea came from the example.}
\end{itemize}
\end{safety}

\section{Guided demo (10--13 min)}
\begin{workedexample}[Model the 2-minute demo]
Run the Night Sky Postcard (in \texttt{project/reference/}) and follow the script, so students hear what \enquote{good} sounds like:
\begin{enumerate}[leftmargin=1.4em,itemsep=1pt]
  \item \enquote{Hi! My project is called Night Sky Postcard. It draws houses under a sky full of stars.}
  \item \enquote{It will ask you two questions. Do you want to answer them?} (A student answers \texttt{25} and \texttt{night}.)
  \item \enquote{Because you said night, it drew a black sky and a moon. That's my \texttt{if}.}
  \item Point at the code: \enquote{This is my function \texttt{house}. It draws one house. Here I use it four times.}
  \item \enquote{One bug I fixed: my stars were all in one place until I used \texttt{random}.} \enquote{Any questions?}
\end{enumerate}
\end{workedexample}

\section{Hands-on build: showcase (13--35 min)}
\begin{unplugged}[Gallery walk with feedback cards]
Visitors (other students and family) walk from stand to stand. At each stand: the artist gives the 2-minute demo, the visitor answers the program's questions, asks one question, and fills in a feedback card (two stars and a wish, signed with a first name or \enquote{a visitor}). Ring a bell or chime every 4 minutes to move on. With no family present, students take turns: half are artists, half visitors, then swap; with 1 student, you (and any family member) are the visitors.
\end{unplugged}
While they present, use the rubric in the Project Teacher Guide (Explaining row) and tick M4 on each brief. Make sure every artist gets the same attention and at least two feedback cards.

\section{Exit assessment (35--55 min)}
Hand out the \textbf{Exit assessment in the Assessments booklet}. Families can chat or leave; students work alone and quietly for about 20 minutes. Read each question aloud if a student asks --- the assessment checks coding, not reading. Students may use Thonny only where the paper says so. Mark later with the Answer Key; it does not change the certificate.
\begin{teachernote}[Keep it low-stakes]
Say: \enquote{This is not a test you pass or fail. It shows me what you've learned and what we should practise in Level 2. Do what you can; skip what you don't know yet and come back.}
\end{teachernote}

\section{Wrap-up (55--60 min)}
Read each artist's feedback cards aloud (or hand them over), then present the certificates one by one with a sentence about something specific they did (\enquote{Kai, your fish loop made the whole sea move}). Everyone claps.
\begin{formativecheck}[Exit ticket: one round]
Each student says: \enquote{One thing I learned in Level 1 is \ldots{} One thing I want to make next is \ldots} Note the answers: they guide Level 2 and the holiday challenge menu.
\end{formativecheck}
\begin{reflection}
\begin{itemize}[leftmargin=1.2em,itemsep=1pt]
  \item Look back at your week 1 robot drawing. What can you do now that you couldn't then?
  \item Which bug are you proudest of fixing?
\end{itemize}
\textbf{Level 2 preview} (one minute): \enquote{In Level 2 you'll make your own arcade games and art studios: lists of things, games that react to keys and timers, programs that save high scores to files --- and you'll even write tests like the \texttt{check\_my\_work} checks.} Hand out homework 12 (reflection + holiday \enquote{keep coding} menu).
\end{reflection}

\section{Notes for the teacher}
\begin{commonmistakes}
  \item \textbf{Last-minute big changes} that break the project. Fix: the 3-minute polish rule; if it breaks, open the last saved version.
  \item \textbf{Demo becomes a code reading of every line.} Fix: \enquote{Pick one part, tell its story.}
  \item \textbf{Turtle window hidden} behind Thonny during the demo: click its icon in the Dock; close it before the next visitor runs it again.
  \item \textbf{Shy presenters freeze.} Let the assistant run the program and ask the artist the script questions as an interview.
\end{commonmistakes}
\begin{differentiation}
\supportpath A student may present in Scratch (support path) or present with a partner, each explaining one part. Shorten the script to three lines: name, what it draws, one part of the code. In the assessment, read every question aloud and accept oral answers.
\extensionpath Ask expert presenters to explain a bug they fixed and why they chose their design. Offer the assessment's stretch question, and suggest a Level 2 starter: \enquote{make your art show redraw when you press the space bar}.
\end{differentiation}
\begin{stretchchallenge}[Visitor's choice]
At the end, let a visitor ask for a tiny change (\enquote{make it purple}, \enquote{100 stars}) and the artist makes it live while explaining what they change.
\end{stretchchallenge}
\begin{teachernote}[If you are short on time]
Run the gallery with 2-minute demos only (no rotation; everyone watches each demo) and move the exit assessment to the start of Buffer B2 or Level 2's first session. Never skip certificates and the round of \enquote{one thing I learned} --- that is the celebration.
\end{teachernote}
